\documentclass[11pt]{article}

\usepackage[a4paper,margin=1.05in]{geometry}
\usepackage{amsmath,amssymb,amsthm,mathtools}
\usepackage{enumitem}
\usepackage{xcolor}
\usepackage[colorlinks=true,linkcolor=blue!50!black,citecolor=blue!50!black,urlcolor=blue!50!black]{hyperref}
\usepackage[T1]{fontenc}
\usepackage[utf8]{inputenc}
\usepackage{lmodern}
\usepackage{microtype}

\newcommand{\R}{\mathbb{R}}
\newcommand{\Rp}{\mathbb{R}_{+}}
\newcommand{\PSD}{\mathsf{PSD}}
\newcommand{\NN}{\mathsf{NN}}
\newcommand{\DNN}{\mathsf{DNN}}
\newcommand{\CP}{\mathsf{CP}}
\newcommand{\DP}{\mathsf{DP}}
\newcommand{\TP}{\mathsf{TP}}
\newcommand{\diag}{\operatorname{diag}}
\newcommand{\Adj}{\operatorname{Adj}}
\newcommand{\one}{\mathbf{1}}

\theoremstyle{definition}
\newtheorem{definition}{Definition}
\newtheorem{example}{Example}
\theoremstyle{plain}
\newtheorem{proposition}{Proposition}
\newtheorem{lemma}{Lemma}
\theoremstyle{remark}
\newtheorem{remark}{Remark}

\title{Doubly Positive and Totally Positive Kernels\\
\large A compact account for positive sketching}
\author{}
\date{}

\begin{document}
\maketitle

\noindent\textbf{Terminology.}
In the matrix literature, the standard names are \emph{doubly nonnegative} and
\emph{completely positive}.  Below I use \emph{doubly positive} and
\emph{totally positive} for the analogous kernel notions, because the latter is
what is needed for positive feature sketching.  To avoid ambiguity, note that
``totally positive'' here does \emph{not} mean Karlin's total positivity, i.e.
nonnegativity of all minors of a two-variable kernel; it means ``completely
positive type'' in the sense of the completely positive matrix cone
\cite{Karlin1968,DeCorte2022}.

\section{The finite-dimensional model}

For a symmetric matrix $A\in\R^{n\times n}$, define
\[
\DNN_n := \{A\in\R^{n\times n}: A\succeq 0,\ A_{ij}\ge 0\ \forall i,j\}
\]
and
\[
\CP_n := \{A\in\R^{n\times n}: A=BB^\top\text{ for some }B\ge 0\}.
\]
Equivalently,
\[
\CP_n=\operatorname{cone}\{vv^\top: v\in\Rp^n\}.
\]
Thus every completely positive matrix is doubly nonnegative:
\[
\CP_n\subseteq \DNN_n.
\]
The classical finite-dimensional fact is
\[
\boxed{\CP_n=\DNN_n\quad\text{for }n\le 4,}
\qquad
\boxed{\CP_n\subsetneq\DNN_n\quad\text{for }n\ge 5.}
\]
See, for example, the monograph of Berman and Shaked-Monderer
\cite{BermanShakedMonderer2003}; the $5\times5$ difference is discussed in
Anstreicher--Burer--Duer \cite{AnstreicherBurerDuer2009}.

\section{Kernel analogues}

Let $X$ be a set and let
\[
k:X\times X\to\R
\]
be symmetric.  For points $x_1,\dots,x_n\in X$, write
\[
K[x_1,\dots,x_n] := \big(k(x_i,x_j)\big)_{i,j=1}^n.
\]

\begin{definition}[Doubly positive kernel]
A kernel $k$ is \emph{doubly positive}, written $k\in\DP(X)$, if
\[
K[x_1,\dots,x_n]\in\DNN_n
\]
for every finite choice $x_1,\dots,x_n\in X$.  Equivalently,
$k$ is positive semidefinite and pointwise nonnegative:
\[
\sum_{i,j=1}^n c_i c_j k(x_i,x_j)\ge 0
\quad\text{for all }c_i\in\R,
\qquad
k(x,y)\ge 0\quad\forall x,y\in X.
\]
\end{definition}

\begin{definition}[Totally positive / completely positive-type kernel]
A kernel $k$ is \emph{totally positive}, written $k\in\TP(X)$, if
\[
K[x_1,\dots,x_n]\in\CP_n
\]
for every finite choice $x_1,\dots,x_n\in X$.
\end{definition}

The inclusion
\[
\TP(X)\subseteq\DP(X)
\]
is immediate from $\CP_n\subseteq\DNN_n$.  Consequently, on a fixed sample of
$n\le4$ points the doubly positive and totally positive tests coincide, whereas
from $n=5$ onward they can differ.  The point is that $\TP(X)$ is the
kernel-level condition corresponding to nonnegative Gram factorizations on every
finite sample.

\section{Why total positivity is the sketching condition}

A positive sketch of $k$ is an integral representation
\begin{equation}\label{eq:positive-sketch}
 k(x,y)
 =
 \int_Z \langle \phi(x,z),\phi(y,z)\rangle_{\R^p}\,d\mu(z),
 \qquad
 \phi(x,z)\in\Rp^p.
\end{equation}
This is the structure used by positive random features: the feature map takes
values in the positive orthant, and finite random features approximate the
integral by a finite sum \cite{ScetbonCuturi2020}.

Fix $x_1,\dots,x_n\in X$.  For each coordinate $a=1,\dots,p$ and latent point
$z\in Z$, set
\[
v_{a,z}:=\big(\phi_a(x_1,z),\dots,\phi_a(x_n,z)\big)^\top\in\Rp^n.
\]
Then \eqref{eq:positive-sketch} gives
\[
K[x_1,\dots,x_n]
=
\sum_{a=1}^p\int_Z v_{a,z}v_{a,z}^\top\,d\mu(z).
\]
Since $\CP_n$ is a closed convex cone, the matrix belongs to $\CP_n$.
Therefore
\[
\boxed{\text{positive sketch} \quad\Longrightarrow\quad k\in\TP(X).}
\]
Conversely, when $X$ itself is finite, this condition is exactly sufficient:
$K\in\CP_n$ means $K=BB^\top$ with $B\ge0$, so one can take the columns of $B$
as positive features.  For general spaces, $\TP(X)$ is the finite-restriction
condition; analytic treatments often define the corresponding cone as the
closure of conic combinations of rank-one kernels $f\otimes f$ with $f\ge0$
\cite{DeCorte2022}.

\section{Complexity}

For a fixed finite sample, checking $K\in\DNN_n$ is tractable: one checks
entrywise nonnegativity and positive semidefiniteness.  Checking
$K\in\CP_n$ is very different.  As $n$ varies, the strong membership problem
for the completely positive cone is NP-hard, and even weak membership for the
completely positive cone and its dual copositive cone is NP-hard
\cite{DickinsonGijben2014}.  Thus, for finite data, membership in $\TP$ is
exactly the computationally difficult part; for an infinite kernel, global
membership in $\TP(X)$ asks for this condition on every finite restriction.

\section{Examples}

\begin{example}[Finite positive features]
If $h:X\to\Rp^r$ and
\[
k(x,y)=\langle h(x),h(y)\rangle_{\R^r},
\]
then $k\in\TP(X)$, with the trivial sketch $Z=\{\ast\}$ and
$\phi(x,\ast)=h(x)$.
\end{example}

\begin{example}[Positive autocorrelations]
Let $X=\R^d$ and $k(x,y)=\psi(x-y)$.  Suppose there are functions
$g_\ell\in L^2(\R^d)$, $g_\ell\ge0$, such that
\[
\psi(t)=\sum_\ell \int_{\R^d} g_\ell(u)g_\ell(u+t)\,du.
\]
Then $k\in\TP(\R^d)$ because
\[
k(x,y)=\sum_\ell \int_{\R^d} g_\ell(u-x)g_\ell(u-y)\,du,
\]
so one may take $z=(u,\ell)$ and
\[
\phi(x,(u,\ell))=g_\ell(u-x)e_\ell.
\]
For instance, $g=\mathbf{1}_{[0,1]}$ on $\R$ gives the triangular kernel
\[
k(x,y)=(1-|x-y|)_+.
\]
\end{example}

\begin{example}[Gaussian kernel]
For
\[
k_\varepsilon(x,y)=\exp\!\left(-\frac{\|x-y\|^2}{2\varepsilon}\right),
\]
set
\[
g_\varepsilon(u)=\left(\frac{2}{\pi\varepsilon}\right)^{d/4}
\exp\!\left(-\frac{\|u\|^2}{\varepsilon}\right).
\]
Then
\[
\int_{\R^d} g_\varepsilon(u-x)g_\varepsilon(u-y)\,du
=
\exp\!\left(-\frac{\|x-y\|^2}{2\varepsilon}\right),
\]
and hence $k_\varepsilon\in\TP(\R^d)$.  This is the clean Lebesgue-measure
version of the Gaussian positive-feature construction.
\end{example}

\begin{example}[Gaussian scale mixtures]
If
\[
k(x,y)=F(\|x-y\|^2),
\qquad
F(s)=\int_0^\infty e^{-as}\,d\nu(a),
\]
for a finite positive measure $\nu$, then $k\in\TP(\R^d)$ by mixing the Gaussian
sketches.  Equivalently, any completely monotone radial kernel
$F(\|x-y\|^2)$ is totally positive.  This includes squared-exponential kernels,
rational-quadratic kernels, inverse-multiquadrics, and Matern kernels via their
standard Gaussian-mixture representations; see, e.g., Rasmussen--Williams for
standard Gaussian-process covariance families \cite{RasmussenWilliams2006}.
\end{example}

\section{Closure properties}

\begin{proposition}[Sums and pointwise products]
Both $\DP(X)$ and $\TP(X)$ are stable under nonnegative sums and pointwise
products.  More precisely, if $k_1,k_2\in\TP(X)$ and $a_1,a_2\ge0$, then
\[
a_1k_1+a_2k_2\in\TP(X),
\qquad
k_1k_2\in\TP(X),
\]
where $(k_1k_2)(x,y)=k_1(x,y)k_2(x,y)$.
\end{proposition}

\begin{proof}
For sums, take the disjoint union of the latent spaces and scale the features
by $\sqrt{a_1}$ and $\sqrt{a_2}$.  For products, if
\[
k_i(x,y)=\int_{Z_i}\langle \phi_i(x,z_i),\phi_i(y,z_i)\rangle\,d\mu_i(z_i),
\]
then
\[
k_1(x,y)k_2(x,y)
=
\int_{Z_1\times Z_2}
\big\langle
\phi_1(x,z_1)\otimes\phi_2(x,z_2),
\phi_1(y,z_1)\otimes\phi_2(y,z_2)
\big\rangle
\,d\mu_1(z_1)d\mu_2(z_2),
\]
and the tensor-product feature is still entrywise nonnegative.  The same
statements for $\DP(X)$ follow from the Schur product theorem and entrywise
nonnegativity.
\end{proof}

\section{Counterexamples}

\subsection*{A finite kernel: doubly positive but not totally positive}
Let $C_5$ be the cycle graph on five vertices and let
\[
A_t=I_5+t\,\Adj(C_5).
\]
The eigenvalues of $\Adj(C_5)$ are $2\cos(2\pi j/5)$, $j=0,\dots,4$; hence
$A_t\succeq0$ for $0\le t\le 1/\varphi$, where
$\varphi=(1+\sqrt5)/2$.  Thus, for example,
\[
A_{3/5}\in\DNN_5.
\]
But $A_t\notin\CP_5$ for $t>1/2$.  Indeed, if $A_t=BB^\top$ with $B\ge0$,
then the zero entries of $A_t$ force every column of $B$ to be supported on a
clique of $C_5$.  Since $C_5$ has no triangles, each column is supported on a
single vertex or on one edge.  To produce the off-diagonal value $t$ on a fixed
edge, the corresponding columns contribute at least $2t$ to the two diagonal
entries, by $u^2+v^2\ge 2uv$.  Summing over the five edges gives
\[
\operatorname{tr}(A_t)\ge 10t.
\]
But $\operatorname{tr}(A_t)=5$, so complete positivity would imply
$t\le1/2$.  Therefore $A_{3/5}$ is doubly nonnegative but not completely
positive.  On the finite space $X=\{1,\dots,5\}$, the kernel
$k(i,j)=(A_{3/5})_{ij}$ is doubly positive but not totally positive.

\subsection*{A translation-invariant counterexample on \texorpdfstring{$\R$}{R}}
The gap is not merely a finite-space pathology.  Define
\[
\psi(t)=\frac{9}{10}\left[
\frac{\sqrt5}{5}
+\frac{5-\sqrt5}{5}\cos\!\left(\frac{2\pi t}{5}\right)
+\frac19\right],
\qquad
k(x,y)=\psi(x-y).
\]
This kernel is continuous, translation-invariant, positive definite, and
strictly positive pointwise: it is a positive linear combination of
$1,e^{2\pi i t/5},e^{-2\pi i t/5}$, and its minimum is positive.  Also
$\psi(0)=1$.

For $x_j=j$, $j=0,\dots,4$, the Gram matrix is
\[
K[x_0,\dots,x_4]=\frac{9}{10}\left(W+\frac19J\right),
\]
where $J$ is the all-ones matrix and
\[
W=
\begin{pmatrix}
1&a&0&0&a\\
a&1&a&0&0\\
0&a&1&a&0\\
0&0&a&1&a\\
a&0&0&a&1
\end{pmatrix},
\qquad
 a=\frac{\sqrt5-1}{2}.
\]
Let
\[
H=
\begin{pmatrix}
1&-1&1&1&-1\\
-1&1&-1&1&1\\
1&-1&1&-1&1\\
1&1&-1&1&-1\\
-1&1&1&-1&1
\end{pmatrix}
\]
be the Horn copositive matrix.  Every completely positive matrix $B$ satisfies
$\langle H,B\rangle\ge0$, while
\[
\left\langle H,\frac{9}{10}\left(W+\frac19J\right)\right\rangle
=\frac{9}{10}\left(5(2-\sqrt5)+\frac59\right)<0.
\]
Hence this finite restriction is not completely positive, and $k$ is doubly
positive but not totally positive.  Thus translation invariance and strict
pointwise positivity do not by themselves imply positive sketchability.

\bibliographystyle{alpha}
\begin{thebibliography}{99}

\bibitem[ABD09]{AnstreicherBurerDuer2009}
Kurt M. Anstreicher, Samuel Burer, and Mirjam Duer.
\newblock The difference between $5\times5$ doubly nonnegative and completely positive matrices.
\newblock \emph{Linear Algebra and its Applications}, 431(9):1539--1552, 2009.

\bibitem[BSM03]{BermanShakedMonderer2003}
Abraham Berman and Naomi Shaked-Monderer.
\newblock \emph{Completely Positive Matrices}.
\newblock World Scientific, 2003.

\bibitem[DGL21]{DeCorte2022}
Erwin DeCorte, Fernando M\'ario de Oliveira Filho, and Frank Vallentin.
\newblock Complete positivity and distance-avoiding sets.
\newblock \emph{Mathematical Programming}, 191:487--558, 2022.

\bibitem[DG14]{DickinsonGijben2014}
Peter J. C. Dickinson and Luuk Gijben.
\newblock On the computational complexity of membership problems for the completely positive cone and its dual.
\newblock \emph{Computational Optimization and Applications}, 57(2):403--415, 2014.

\bibitem[Kar68]{Karlin1968}
Samuel Karlin.
\newblock \emph{Total Positivity, Volume I}.
\newblock Stanford University Press, 1968.

\bibitem[RW06]{RasmussenWilliams2006}
Carl Edward Rasmussen and Christopher K. I. Williams.
\newblock \emph{Gaussian Processes for Machine Learning}.
\newblock MIT Press, 2006.

\bibitem[SC20]{ScetbonCuturi2020}
Meyer Scetbon and Marco Cuturi.
\newblock Linear time Sinkhorn divergences using positive features.
\newblock In \emph{Advances in Neural Information Processing Systems}, 33, 2020.

\end{thebibliography}

\end{document}
